\textbf{Задача 24.} Сформулируйте и докажите теорему о независимости класса $\rp$ от точного значения пороговой вероятности.

\begin{proof}
Пусть $c$ и $d$ целые положительные числа. Тогда
\[\bpp_{0,\frac{1}{n^c}}=\bpp_{0,1 - \frac{1}{2^{n^d}}}=\rp\]
Очевидно, что

\[\bpp_{0,1 - \frac{1}{2^{n^d}}}\subset\rp\subset\bpp_{0,\frac{1}{n^c}}\]

Достаточно показать, что $\bpp_{0,\frac{1}{n^c}}=\bpp_{0,1 - \frac{1}{2^{n^d}}}$

Очевидно, что $\bpp_{0,1 - \frac{1}{2^{n^d}}}\subset\bpp_{0,\frac{1}{n^c}}$

Пусть язык $A\in\bpp_{0,\frac{1}{n^c}}$.
Возьмем машину, которая его распознает, и запустим её $K=n^{c+d}$ раз, результат исполнения положим как дизъюнкция результатов запусков. Тогда вероятность вернуть $1$ для слова, не лежащего в $A$, это $0$. А вероятность вернуть $0$ при слове, лежащем в $A$, меньше или равна
\[
\left(1-\frac{1}{n^c}\right)^K = \exp\left(K\cdot \ln{\left(1 - \frac{1}{n^c}\right)}\right) \le \exp\left(-\frac{K}{n^c}\right) = e^{-n^d} \le 2^{-n^d}
\]

Значит, $\forall x \in A\ \mathrm{Pr}(M(x)=1) \ge 1 - \dfrac{1}{2^{n^d}}$

Значит, $A\in\bpp_{0,1 - \frac{1}{2^{n^d}}}$. Тогда $\bpp_{0,\frac{1}{n^c}}=\bpp_{0,1 - \frac{1}{2^{n^d}}}$. \end{proof}